\documentclass[12pt, a4paper]{article}
\usepackage[utf8]{inputenc}
\usepackage[T1]{fontenc}
\usepackage[romanian]{babel}
\usepackage{amsmath}
\usepackage{amssymb}
\usepackage{geometry}
\geometry{a4paper, margin=2cm}
\usepackage{listings}
\usepackage{xcolor}

\lstset{
    language=Python,
    basicstyle=\ttfamily\small,
    keywordstyle=\color{blue},
    stringstyle=\color{red},
    commentstyle=\color{green!50!black},
    showstringspaces=false,
    breaklines=true,
    frame=single,
    numbers=left,
    numberstyle=\tiny\color{gray}
}

\title{Rezolvare Teme - Algoritmi de Primalitate \c{s}i Factorizare}
\author{}
\date{}

\begin{document}
\maketitle  

\section*{Tema 2: Implementarea algoritmului de factorizare Fermat}

Acest algoritm se bazeaz\u{a} pe g\u{a}sirea unei scrieri a lui $n$ ca diferen\c{t}\u{a} de p\u{a}trate: $n = t^2 - s^2 = (t-s)(t+s)$. C\u{a}utarea \^{i}ncepe de la valoarea $t = \lceil\sqrt{n}\rceil$.

\begin{lstlisting}
import math

def factorizare_fermat(n):
    if n % 2 == 0: return (n // 2, 2)
    t = math.isqrt(n)
    if t * t < n:
        t += 1
        
    while True:
        s2 = t**2 - n
        s = math.isqrt(s2)
        if s**2 == s2:
            a = t - s
            b = t + s
            return (int(a), int(b))
        t += 1
\end{lstlisting}

\end{document}